\subsection{The compressed dynamics are nonlinear and feature-learning}
\label{sec:nonlinear-learning}

A moving kernel alone does not distinguish nonlinear learning from a deep
linear network. We therefore separate two comparisons. In this subsection,
$f_{\mathrm{NTK},n}$ is the predictor obtained by freezing the coupled
dense reference's initial tangent kernel. Since $w(0)=0$, this is exactly its readout-only flow
with all hidden features fixed at initialization.

\begin{samepage}
\begin{theorem}[Compression preserves nonlinear feature learning]
\label{thm:nonlinear-learning}
Under the compression theorem's assumptions, suppose additionally that every
activation is nonaffine and $|x_a^\top x_b|<d$ for $a\ne b$.
For Taylor, augment $\mathcal X$ by at most four deterministic passive
witness inputs, declared before initialization without their labels.
There are $c_*,t_*>0$, depending on the fixed problem but not on width,
such that for every fixed $0<t\le t_*$, each compressed predictor satisfies,
with probability tending to one as $n\to\infty$,
\begin{align}
 \inf_{a\in\R^d,\,b\in\R}\sup_{x\in\mathcal X}
 |f_{\rm model}(t,x)-a^\top x-b|&\ge c_*t,
 \label{fl:input-gap}\\
 \sup_{x\in\mathcal X}
 |f_{\rm model}(t,x)-f_{\mathrm{NTK},n}(t,x)|&\ge c_*t^3.
 \label{fl:kernel-gap}
\end{align}
The storage orders and approximation conclusions are unchanged; Taylor's
bound uses the augmented panel count, at most $p+4$ instead of $p$.

In the dense reference, every hidden layer satisfies
\begin{equation}
 \left(\frac1n\sum_{a=1}^m
 \|h^{(\ell)}(t,x_a)-h^{(\ell)}(0,x_a)\|_2^2\right)^{1/2}
 \ge c_*t^2\qquad(1\le\ell\le L)
 \label{fl:layer-motion}
\end{equation}
with the same probability convention. In the first layer, the nonlinear
part of the feature increment itself has positive mean-square size of order
$t^4$, in the precise sense of \eqref{fl:nonlinear-increment}.
\end{theorem}
\end{samepage}

The first inequality excludes every deep linear predictor, even with biases,
an arbitrary optimizer and time-dependent weights. The second excludes the
reference network's frozen-feature training. The internal statement further
rules out a frozen nonlinear feature component accompanied only by an affine
feature change. It does not identify the compressed coordinates with dense
neurons. The additional hypotheses belong only to this activity theorem;
the broader compression theorem does not require them.

The mechanism has three time scales. The initially zero readout first grows
at order $t$; hidden features move at order $t^2$; the output departs from
frozen-feature training at order $t^3$. The leading departure in the label
direction is a sum of squared hidden-layer accelerations, so contributions
from different layers cannot cancel. Separately, the initial feature kernel
has nonlinear input components that no nonzero label vector can cancel
completely on the sphere. Both signals persist while compression error
vanishes. The proof in \Cref{sec:nonlinear-learning-proof} uses only
initial Gaussian conditioning and real short-time estimates, not a
width-independent complex-time radius or a trained population-flow theorem.
These are early-time certificates, not endpoint or test-risk comparisons.
